\documentclass[11pt]{article}

\usepackage[margin=1in]{geometry}
\usepackage{amsmath,amssymb,mathtools,bm}
\usepackage{booktabs}
\usepackage{xcolor}
\usepackage{hyperref}
\usepackage{listings}

\hypersetup{
  colorlinks=true,
  linkcolor=blue!50!black,
  citecolor=blue!50!black,
  urlcolor=blue!50!black,
  pdftitle={Molecular Response Notes},
}

\numberwithin{equation}{section}
\setlength{\parindent}{0pt}
\setlength{\parskip}{0.6em}

% Code style for snippets, file names, and identifiers.
\lstdefinestyle{respcode}{
  basicstyle=\ttfamily\small,
  keywordstyle=\color{blue!60!black},
  commentstyle=\color{green!40!black},
  stringstyle=\color{red!50!black},
  numbers=left,
  numberstyle=\tiny\color{gray},
  numbersep=8pt,
  breaklines=true,
  showstringspaces=false,
  frame=single,
  columns=fullflexible,
  tabsize=2
}

% Text helpers for code-facing notes.
\newcommand{\code}[1]{\texttt{#1}}
\newcommand{\class}[1]{\texttt{#1}}
\newcommand{\func}[1]{\texttt{#1}}
\newcommand{\file}[1]{\texttt{#1}}
\newcommand{\param}[1]{\texttt{#1}}

% General math helpers.
\newcommand{\dd}{\mathop{}\!\mathrm{d}}
\newcommand{\ii}{\mathrm{i}}
\newcommand{\ee}{\mathrm{e}}
\newcommand{\trans}{^{\mathsf{T}}}
\newcommand{\herm}{^{\dagger}}
\newcommand{\inv}{^{-1}}
\newcommand{\vect}[1]{\bm{#1}}
\newcommand{\mat}[1]{\bm{#1}}
\newcommand{\op}[1]{\hat{#1}}
\newcommand{\diag}{\operatorname{diag}}
\newcommand{\Tr}{\operatorname{Tr}}
\newcommand{\Repart}{\operatorname{Re}}
\newcommand{\Impart}{\operatorname{Im}}

% Quantum chemistry notation.
\newcommand{\bra}[1]{\langle #1 \rvert}
\newcommand{\ket}[1]{\lvert #1 \rangle}
\newcommand{\braket}[2]{\langle #1 \mid #2 \rangle}
\newcommand{\expect}[1]{\langle #1 \rangle}
\newcommand{\ao}[1]{\chi_{#1}}
\newcommand{\mo}[1]{\phi_{#1}}
\newcommand{\state}[1]{\ket{\Psi_{#1}}}
\newcommand{\gs}{\Psi_0}
\newcommand{\occset}{\mathcal{O}}
\newcommand{\virtset}{\mathcal{V}}
\newcommand{\basisset}{\mathcal{B}}
\newcommand{\nocc}{N_{\mathrm{occ}}}
\newcommand{\nvirt}{N_{\mathrm{virt}}}
\newcommand{\nmo}{N_{\mathrm{MO}}}

% Common matrices and vectors in response theory.
\newcommand{\ident}{\mat{I}}
\newcommand{\coef}{\mat{C}}
\newcommand{\density}{\mat{D}}
\newcommand{\fock}{\mat{F}}
\newcommand{\overlap}{\mat{S}}
\newcommand{\coreH}{\mat{H}}
\newcommand{\coulomb}{\mat{J}}
\newcommand{\exchange}{\mat{K}}
\newcommand{\Aresp}{\mat{A}}
\newcommand{\Bresp}{\mat{B}}
\newcommand{\Mresp}{\mat{M}}
\newcommand{\metric}{\mat{\Sigma}}
\newcommand{\trial}{\mat{V}}
\newcommand{\subspace}{\mat{W}}
\newcommand{\Xresp}{\vect{X}}
\newcommand{\Yresp}{\vect{Y}}
\newcommand{\rhs}{\vect{b}}
\newcommand{\residual}{\vect{r}}

% Response-specific shorthand.
\newcommand{\freq}{\omega}
\newcommand{\freqshift}{\eta}
\newcommand{\exenergy}[1]{\Omega_{#1}}
\newcommand{\pert}[1]{\op{V}^{#1}}
\newcommand{\prop}[1]{\alpha_{#1}}
\newcommand{\response}[2]{#1\!\left(#2\right)}
\newcommand{\lrmat}{\Aresp - (\freq + \ii \freqshift)\metric}

\title{Molecular Response Notes}
\author{Your Name}
\date{\today}

\begin{document}

\maketitle
\tableofcontents

\section{Purpose}

These notes are a working document for the molecular response implementation.
Keep theory, algorithm choices, code structure, and validation cases in one place.
The preamble is intentionally macro-heavy so notation changes can be made once and
applied everywhere.

\section{Notation Cheat Sheet}

\begin{table}[h]
\centering
\begin{tabular}{ll}
\toprule
Notation & Meaning \\
\midrule
$\ao{\mu}, \ao{\nu}$ & atomic-orbital basis functions \\
$\mo{p}, \mo{q}$ & molecular orbitals \\
$\coef$ & AO-to-MO coefficient matrix \\
$\density$ & density matrix \\
$\fock$ & effective one-body / Fock-like matrix \\
$\Aresp, \Bresp$ & response blocks \\
$\Xresp, \Yresp$ & excitation or response amplitudes \\
$\freq$ & perturbation frequency \\
\code{\class{ResponseVector}} & code identifier recorded with a text macro \\
\bottomrule
\end{tabular}
\end{table}

\section{Core Equations}

\subsection{Orbital Expansion}

\begin{equation}
  \mo{p}(\vect{r}) = \sum_{\mu \in \basisset} C_{\mu p}\,\ao{\mu}(\vect{r}).
\end{equation}

\subsection{Frequency-Domain Linear Response}

\begin{equation}
  \left[\Aresp - (\freq + \ii \freqshift)\metric\right]
  \response{\Xresp}{\freq}
  =
  \response{\rhs}{\freq}.
  \label{eq:linear-response}
\end{equation}

Document here exactly what \(\Aresp\), \(\metric\), and \(\rhs\) mean in your
implementation. If the code uses a reduced space, this is also the right place
to define the projected problem.

\subsection{Excited-State Eigenproblem}

\begin{equation}
  \begin{pmatrix}
    \Aresp & \Bresp \\
    \Bresp^{*} & \Aresp^{*}
  \end{pmatrix}
  \begin{pmatrix}
    \Xresp_{n} \\
    \Yresp_{n}
  \end{pmatrix}
  =
  \exenergy{n}
  \begin{pmatrix}
    \ident & 0 \\
    0 & -\ident
  \end{pmatrix}
  \begin{pmatrix}
    \Xresp_{n} \\
    \Yresp_{n}
  \end{pmatrix}.
  \label{eq:excited-state}
\end{equation}

\subsection{Transition Property Placeholder}

\begin{equation}
  \vect{\mu}_{n}
  =
  \bra{\gs}\op{\vect{\mu}}\state{n}.
  \label{eq:transition-dipole}
\end{equation}

Replace Eq.~\eqref{eq:transition-dipole} with the property expression that is
most useful for your current work: polarizability, oscillator strengths, Raman
intensities, or response residues.

\section{Implementation Map}

\subsection{Key Objects}

\begin{itemize}
  \item \class{ResponseVector}: note storage layout, ownership, and normalization.
  \item \class{ExcitedStateBundleSolver}: note subspace growth, restart, and locking.
  \item \func{apply\_response\_operator(...)}: describe the mathematical operator
        applied by the code.
  \item \file{src/apps/molresponse\_v2/...}: keep exact file paths once they settle.
\end{itemize}

\subsection{Code Entry Template}

\begin{lstlisting}[style=respcode,language=C++,caption={Starter block for recording an implementation detail}]
// File: src/apps/molresponse_v2/ExcitedStateBundleSolver.cpp
// Routine: solve_bundle(...)
// Role: build the projected problem and update bundle residuals
// Inputs: trial subspace, operator application, convergence thresholds
// Outputs: Ritz pairs, residual norms, restart metadata
\end{lstlisting}

\section{Algorithm Notes}

\subsection{Data Flow}

\begin{enumerate}
  \item Build or update the trial space \(\trial\).
  \item Apply the operator to form projected quantities.
  \item Solve the reduced problem.
  \item Form \(\residual\) and decide expand / lock / restart.
\end{enumerate}

\subsection{Numerical Questions}

\begin{itemize}
  \item What is the effective preconditioner for the current operator?
  \item Which quantities must be orthogonalized across frequencies or states?
  \item Where does symmetry reduce the cost or storage footprint?
  \item What is the failure mode when the projected metric becomes ill-conditioned?
\end{itemize}

\section{Validation Plan}

\begin{itemize}
  \item Small molecules with reference excitation energies.
  \item Frequency-domain checks against finite differences.
  \item Restart and checkpoint consistency.
  \item Invariance tests under subspace reordering or block partition changes.
\end{itemize}

\section{Open Items}

\begin{itemize}
  \item \textbf{TODO:} record the exact response equations used in the current code path.
  \item \textbf{TODO:} document how the solver labels frequencies, states, and bundles.
  \item \textbf{TODO:} add the file/function map once interfaces stop moving.
\end{itemize}

\appendix

\section{Scratch Derivations}

Use this appendix for algebra that is not ready for the main narrative.
Once a derivation stabilizes, move the cleaned version into the main sections.

\end{document}
